\SourceLeaf{464}
\par\Needspace{8\baselineskip}\bigskip
\pdfbookmark[3]{Darstellung, 15}{pallas-II-article15}
\begin{center}15.\end{center}
\noindent\hspace*{\parindent}Die bisherigen Untersuchungen sind von der Natur der störenden Kräfte unabhängig. Nehmen wir nun an, dass sie durch die Einwirkung eines anderen Körpers sowohl auf den Planeten als auch auf die Sonne entstehen. Es sei \(\mu m\) die Masse dieses Körpers, \(x',y',z'\) seine Koordinaten in Bezug auf die drei Fundamentalebenen, \(r'=\BareRoot{(x'x'+y'y'+z'z')}\) sein Abstand von der Sonne und \(\rho=\BareRoot{((x'-x)^2+(y-y')^2+(z-z')^2)}\) sein Abstand vom gestörten Planeten. Nach bekannten Grundsätzen gilt dann
\[
\begin{aligned}
\xi&=\dfrac{\mu mx'}{r'^3}-\dfrac{\mu m(x'-x)}{\rho^3}\\
\eta&=\dfrac{\mu my'}{r'^3}-\dfrac{\mu m(y'-y)}{\rho^3}\\
\zeta&=\dfrac{\mu mz'}{r'^3}-\dfrac{\mu m(z'-z)}{\rho^3}\mathpunct{\cdot}
\end{aligned}
\]
\par Bezeichnen wir den heliozentrischen Ort des störenden Planeten mit \(L'\), so gilt
\[
\begin{aligned}
x'&=r'\cos XL'\\
y'&=r'\cos YL'\\
z'&=r'\cos ZL';
\end{aligned}
\]
Die Formeln aus Artikel 9 ergeben also unter Berücksichtigung von \(\Lambda L=ML=90^0\)
\[
\begin{aligned}
T&=\mu\left(\dfrac{1}{\rho^3}-\dfrac{1}{r'^3}\right)r'\cos LL'-\dfrac{\mu r}{\rho^3}\\
V&=\mu\left(\dfrac{1}{\rho^3}-\dfrac{1}{r'^3}\right)r'\cos\Lambda L'\\
W&=\mu\left(\dfrac{1}{\rho^3}-\dfrac{1}{r'^3}\right)r'\cos ML'.
\end{aligned}
\]
Es sei \(\beta'\) die Neigung des Radiusvektors \(r'\) gegen die Ebene der oskulierenden Bahn des gestörten Planeten (mit dem Vorzeichen \(+\) auf der Seite des Pols \(M\)). Seine Projektion auf diese Ebene bilde mit der Linie des aufsteigenden Knotens derselben Ebene in der Koordinatenebene \(x,y\) den Winkel \(v'-\AscendingNode\), gemessen in Bewegungsrichtung des gestörten Planeten. Dann lassen sich \(\beta',v'\) als heliozentrische Breite und Länge des störenden Körpers in Bezug \SourceLeaf{465}auf die Ebene der oskulierenden Bahn des gestörten Planeten auffassen. Man sieht leicht, dass
\[
\begin{aligned}
\cos LL'&=\cos\beta'\cos(v'-v)\\
\cos\Lambda L'&=\cos\beta'\sin(v'-v)\\
\cos ML'&=\sin\beta'.
\end{aligned}
\]
\par Um schließlich alles für die Rechnung Erforderliche hier zusammenzustellen, setzen wir die Neigung der Ebene, in der sich der störende Körper bewegt, gegen die Koordinatenebene \(x,y\) gleich \(i'\). Ferner sei \(\AscendingNode'\) die von der \(x\)-Achse aus gemessene Länge des aufsteigenden Knotens. Es sei außerdem \(w'-\AscendingNode'\) der Abstand des heliozentrischen Ortes \(L'\) von diesem Knoten, sodass \(w'\) seine sogenannte wahre Länge in der Bahn ist. Die Größen \(N,\Delta,J\) bestimme man aus den folgenden Gleichungen.
\[
\begin{aligned}
\sin\tfrac12J\mathbin{.}\sin\tfrac12(\Delta+N)&=\sin\tfrac12(i'+i)\mathbin{.}\sin\tfrac12(\AscendingNode'-\AscendingNode)\\
\sin\tfrac12J\mathbin{.}\cos\tfrac12(\Delta+N)&=\sin\tfrac12(i'-i)\mathbin{.}\cos\tfrac12(\AscendingNode'-\AscendingNode)\\
\cos\tfrac12J\mathbin{.}\sin\tfrac12(\Delta-N)&=\cos\tfrac12(i'+i)\mathbin{.}\sin\tfrac12(\AscendingNode'-\AscendingNode)\\
\cos\tfrac12J\mathbin{.}\cos\tfrac12(\Delta-N)&=\cos\tfrac12(i'-i)\mathbin{.}\cos\tfrac12(\AscendingNode'-\AscendingNode).
\end{aligned}
\]
\par Damit erhält man
\[
\begin{aligned}
\operatorname{tang}(v'-\AscendingNode-\Delta)&=\cos J\mathbin{.}\operatorname{tang}(w'-\AscendingNode'-N)\\
\operatorname{tang}\beta'&=\operatorname{tang}J\mathbin{.}\sin(v'-\AscendingNode-\Delta)\\
\sin\beta'&=\sin J\mathbin{.}\sin(w'-\AscendingNode'-N)\\
\rho&=\BareRoot{(rr+r'r'-2rr'\cos\beta'\cos(v'-v))}.
\end{aligned}
\]
